\subsection{Clases}

\subsubsection{Regla: PascalCase y sustantivo singular}

\textbf{Regla:} El nombre de una clase deberá escribirse en \texttt{PascalCase} y ser un sustantivo o frase nominal en singular.

\textbf{Descripción:} Las guías de diseño de .NET establecen que las clases deben nombrarse con sustantivos o frases nominales en \texttt{PascalCase}, distinguiéndolas así de los métodos, que se nombran con frases verbales (Cwalina \& Abrams, 2008).

\textbf{Código correcto:}
\begin{lstlisting}[style=csharpstyle]
public sealed class BoardTile
{
}
\end{lstlisting}

\textbf{Código incorrecto:}
\begin{lstlisting}[style=csharpstyle]
public sealed class BoardTiles
{
}
\end{lstlisting}

\subsubsection{Regla: sin prefijos ni abreviaturas poco claras}

\textbf{Regla:} No se utilizarán prefijos para indicar que un tipo es una clase (por ejemplo, \texttt{CTile} o \texttt{ClsPlayer}), ni abreviaturas que no sean ampliamente reconocidas.

\textbf{Descripción:} Las guías de diseño de .NET indican expresamente que los nombres de clase no deben llevar prefijos que identifiquen su tipo (Cwalina \& Abrams, 2008).

\textbf{Código correcto:}
\begin{lstlisting}[style=csharpstyle]
public sealed class TurnManager
{
}
\end{lstlisting}

\textbf{Código incorrecto:}
\begin{lstlisting}[style=csharpstyle]
public sealed class ClsTrnMgr
{
}
\end{lstlisting}
